\section{Bimodules}

\begin{definition}
    Let $R$ and $S$ be rings. An abelian group $M$ is called a \emph{left-$R$-right-$S$-bimodule} if $M$ is both a left $R$-module and a right $S$-module such that the two scalar multiplications satisfy
    \[
        r(xs)=(rx)s
        \qquad
        \text{for all } r\in R,\ s\in S,\ x\in M.
    \]
\end{definition}

\noindent\textbf{Notation.}
\[
    {}_R M \qquad \text{left } R\text{-module},
\]
\[
    M_S \qquad \text{right } S\text{-module},
\]
\[
    {}_R M_S \qquad \text{left } R\text{-right } S\text{-bimodule}.
\]

Let $R,S$ be rings, and let ${}_R U_S$ be a bimodule.

For each left $R$-module ${}_R M$, there are two naturally induced $S$-module structures:

\[
    \operatorname{Hom}_R({}_R U_S, {}_R M),
    \qquad \text{a left } S\text{-module},
\]

\[
    \operatorname{Hom}_R({}_R M, {}_R U_S),
    \qquad \text{a right } S\text{-module}.
\]

More precisely, for each $s\in S$:

if
\[
    f\in \operatorname{Hom}_R({}_R U_S, {}_R M),
\]
define
\[
    s\cdot f: U\to M,\qquad u\mapsto f(us).
\]

Similarly, if
\[
    f\in \operatorname{Hom}_R({}_R M, {}_R U_S),
\]
define
\[
    f\cdot s: M\to U,\qquad m\mapsto f(m)s.
\]

Thus we obtain functors
\[
    \operatorname{Hom}_R({}_R U_S,-): R\text{-}\mathbf{Mod}\longrightarrow S\text{-}\mathbf{Mod},
\]
given on objects by
\[
    {}_R M\longmapsto \operatorname{Hom}_R({}_R U_S, {}_R M),
\]
which is covariant, and
\[
    \operatorname{Hom}_R(-,{}_R U_S): R\text{-}\mathbf{Mod}\longrightarrow \mathbf{Mod}\text{-}S,
\]
given on objects by
\[
    {}_R M\longmapsto \operatorname{Hom}_R({}_R M, {}_R U_S),
\]
which is contravariant.

On morphisms, for
\[
    f:M\to M',
\]
the induced map
\[
    \operatorname{Hom}_R({}_R U_S, {}_R M)\longrightarrow \operatorname{Hom}_R({}_R U_S, {}_R M')
\]
is given by
\[
    h\longmapsto f\circ h.
\]

the induced map
\[
    \operatorname{Hom}_R({}_R M', {}_R U_S)\longrightarrow \operatorname{Hom}_R({}_R M, {}_R U_S)
\]
is given by
\[
    h\longmapsto  h \circ f.
\]

\begin{theorem}
    The functor
    \[
        \operatorname{Hom}_R({}_R U_S,-): R\text{-}\mathbf{Mod}\longrightarrow S\text{-}\mathbf{Mod}
    \]
    is left exact.
\end{theorem}

\begin{proof}
    Let
    \[
        0\longrightarrow N' \xrightarrow{\,f\,} N \xrightarrow{\,g\,} N''
    \]
    be an exact sequence in $R$-\textbf{Mod}. Consider the induced sequence
    \[
        0\longrightarrow \operatorname{Hom}_R({}_R U_S,N')
        \xrightarrow{\operatorname{Hom}_R(U,f)}
        \operatorname{Hom}_R({}_R U_S,N)
        \xrightarrow{\operatorname{Hom}_R(U,g)}
        \operatorname{Hom}_R({}_R U_S,N'').
    \]
    We show that this sequence is exact.

    First, $\operatorname{Hom}_R(U,f)$ is injective. Let
    \[
        \varphi\in \operatorname{Hom}_R(U,N')
    \]
    and suppose that
    \[
        \operatorname{Hom}_R(U,f)(\varphi)=f\circ \varphi=0.
    \]
    Then for every $u\in U$,
    \[
        f(\varphi(u))=0.
    \]
    Since $f$ is injective, it follows that $\varphi(u)=0$ for all $u\in U$. Hence $\varphi=0$, so $\operatorname{Hom}_R(U,f)$ is injective.

    Next, we show that
    \[
        \operatorname{Im}\bigl(\operatorname{Hom}_R(U,f)\bigr)
        \subseteq
        \ker\bigl(\operatorname{Hom}_R(U,g)\bigr).
    \]
    Indeed, if
    \[
        \varphi\in \operatorname{Hom}_R(U,N'),
    \]
    then
    \[
        \operatorname{Hom}_R(U,g)\bigl(\operatorname{Hom}_R(U,f)(\varphi)\bigr)
        =
        g\circ f\circ \varphi
        =0,
    \]
    because $g\circ f=0$. Therefore
    \[
        \operatorname{Im}\bigl(\operatorname{Hom}_R(U,f)\bigr)
        \subseteq
        \ker\bigl(\operatorname{Hom}_R(U,g)\bigr).
    \]

    Finally, let
    \[
        \psi\in \ker\bigl(\operatorname{Hom}_R(U,g)\bigr).
    \]
    Then $\psi\in \operatorname{Hom}_R(U,N)$ and
    \[
        g\circ \psi=0.
    \]
    Thus for every $u\in U$,
    \[
        g(\psi(u))=0,
    \]
    so $\psi(u)\in \ker g=\operatorname{Im} f$ by exactness. Hence for each $u\in U$ there exists some $y\in N'$ such that
    \[
        f(y)=\psi(u).
    \]
    Since $f$ is injective, such $y$ is unique. Define
    \[
        \varphi:U\to N',\qquad u\mapsto y,
    \]
    where $y$ is the unique element satisfying $f(y)=\psi(u)$.

    Then
    \[
        f\circ \varphi=\psi.
    \]
    It remains to check that $\varphi$ is $R$-linear. For $u,v\in U$,
    \[
        f(\varphi(u+v))=\psi(u+v)=\psi(u)+\psi(v)=f(\varphi(u))+f(\varphi(v))
        =f(\varphi(u)+\varphi(v)).
    \]
    Since $f$ is injective,
    \[
        \varphi(u+v)=\varphi(u)+\varphi(v).
    \]
    Similarly, for $r\in R$ and $u\in U$,
    \[
        f(\varphi(ru))=\psi(ru)=r\psi(u)=rf(\varphi(u))=f(r\varphi(u)),
    \]
    so again by injectivity of $f$,
    \[
        \varphi(ru)=r\varphi(u).
    \]
    Thus $\varphi\in \operatorname{Hom}_R(U,N')$, and
    \[
        \operatorname{Hom}_R(U,f)(\varphi)=f\circ \varphi=\psi.
    \]
    Therefore
    \[
        \ker\bigl(\operatorname{Hom}_R(U,g)\bigr)
        \subseteq
        \operatorname{Im}\bigl(\operatorname{Hom}_R(U,f)\bigr).
    \]

    Combining the two inclusions, we get
    \[
        \operatorname{Im}\bigl(\operatorname{Hom}_R(U,f)\bigr)
        =
        \ker\bigl(\operatorname{Hom}_R(U,g)\bigr).
    \]
    Hence
    \[
        0\longrightarrow \operatorname{Hom}_R({}_R U_S,N')
        \xrightarrow{\operatorname{Hom}_R(U,f)}
        \operatorname{Hom}_R({}_R U_S,N)
        \xrightarrow{\operatorname{Hom}_R(U,g)}
        \operatorname{Hom}_R({}_R U_S,N'')
    \]
    is exact, and so the functor $\operatorname{Hom}_R({}_R U_S,-)$ is left exact.
\end{proof}

\begin{theorem}
    Let ${}_R U_S$ be an $R$-$S$-bimodule. Then the contravariant functor
    \[
        \Hom_R(-,U): R\text{-}\mathrm{Mod}\longrightarrow \mathrm{Mod}\text{-}S
    \]
    is left exact.
\end{theorem}

\begin{proof}
    Since $\Hom_R(-,U)$ is contravariant, to say that it is left exact means the following:

    for every short exact sequence of left $R$-modules
    \[
        A \xrightarrow{\,f\,} B \xrightarrow{\,g\,} C  \longrightarrow 0,
    \]
    the induced sequence of right $S$-modules
    \[
        0 \longrightarrow \Hom_R(C,U)
        \xrightarrow{\ \Hom_R(g,U)\ }
        \Hom_R(B,U)
        \xrightarrow{\ \Hom_R(f,U)\ }
        \Hom_R(A,U)
    \]
    is exact.

    We verify this.

    First, $\Hom_R(g,U)$ is injective. Indeed, let $\varphi \in \Hom_R(C,U)$ and suppose
    \[
        \Hom_R(g,U)(\varphi)=\varphi\circ g=0.
    \]
    Since $g$ is surjective, it follows that $\varphi=0$. Hence $\Hom_R(g,U)$ is injective.

    Next, we show
    \[
        \operatorname{Im}\bigl(\Hom_R(g,U)\bigr)\subseteq \ker\bigl(\Hom_R(f,U)\bigr).
    \]
    Let $\varphi\in \Hom_R(C,U)$. Then
    \[
        \Hom_R(f,U)\bigl(\Hom_R(g,U)(\varphi)\bigr)
        =
        (\varphi\circ g)\circ f
        =
        \varphi\circ (g\circ f)
        =
        0,
    \]
    because $g\circ f=0$. Therefore
    \[
        \operatorname{Im}\bigl(\Hom_R(g,U)\bigr)\subseteq \ker\bigl(\Hom_R(f,U)\bigr).
    \]

    Finally, we prove the reverse inclusion. Let $\psi\in \Hom_R(B,U)$ satisfy
    \[
        \Hom_R(f,U)(\psi)=\psi\circ f=0.
    \]
    Then $\psi$ vanishes on $\operatorname{Im}(f)=\ker(g)$. Hence $\psi$ factors through the quotient
    \[
        B/\ker(g).
    \]
    Since $g$ is surjective, there exists a unique $R$-module homomorphism $\overline{\psi}:C\to U$ such that
    \[
        \psi=\overline{\psi}\circ g.
    \]
    Thus
    \[
        \psi\in\operatorname{Im}\bigl(\Hom_R(g,U)\bigr),
    \]
    so
    \[
        \ker\bigl(\Hom_R(f,U)\bigr)\subseteq\operatorname{Im}\bigl(\Hom_R(g,U)\bigr).
    \]

    Combining the two inclusions, we obtain
    \[
        \operatorname{Im}\bigl(\Hom_R(g,U)\bigr)=\ker\bigl(\Hom_R(f,U)\bigr).
    \]
    Hence the sequence
    \[
        0 \longrightarrow \Hom_R(C,U)
        \xrightarrow{\ \Hom_R(g,U)\ }
        \Hom_R(B,U)
        \xrightarrow{\ \Hom_R(f,U)\ }
        \Hom_R(A,U)
    \]
    is exact. This proves that $\Hom_R(-,U)$ is left exact.
\end{proof}

Let ${}_SM_R$ be a left-$S$-right-$R$-bimodule,${}_R N$ be a left $R$-module, consider the tensor product $M \otimes_R N$.

For each \(s \in S\), the mapping
\[
    M \times N \longrightarrow M \otimes_R N,
    \qquad
    (m,n) \longmapsto (sm)\otimes n
\]
is balanced.

So there is a unique abelian group homomorphism
\[
    u(s): M \otimes_R N \longrightarrow M \otimes_R N
\]
such that
\[
    u(s)(m\otimes n)=(sm)\otimes n.
\]

Only need to check:
\[
    u:S\longrightarrow \operatorname{End}_{\mathbb Z}(M\otimes_R N)
\]
is a ring homomorphism.

If \(M_R,\ {}_R N_T\), then \(M\otimes_R N\) is a right \(T\)-module:
\[
    (m\otimes n)t = m\otimes (nt).
\]

\begin{lemma}
    Let ${}_S M_R$ and ${}_R N_T$ be bimodules. Then $M \otimes_R N$ carries a natural structure of a left $S$-right $T$-bimodule, defined on elementary tensors by
    \[
        s(m\otimes n)=(sm)\otimes n,
        \qquad
        (m\otimes n)t=m\otimes (nt),
    \]
    for all $s\in S$,$t\in T$,$m\in M$, and $n\in N$.
\end{lemma}

\begin{proof}
    Fix $s\in S$. Define
    \[
        \beta_s:M\times N\to M\otimes_R N,
        \qquad
        \beta_s(m,n)=sm\otimes n.
    \]
    We claim that $\beta_s$ is $R$-balanced. Indeed, it is additive in each variable, and for $r\in R$,
    \[
        \beta_s(mr,n)=s(mr)\otimes n=(sm)r\otimes n=sm\otimes rn=\beta_s(m,rn).
    \]
    Hence, by the universal property of the tensor product, there exists a unique group homomorphism
    \[
        \lambda_s:M\otimes_R N\to M\otimes_R N
    \]
    such that
    \[
        \lambda_s(m\otimes n)=sm\otimes n.
    \]
    We write $\lambda_s(x)=sx$. Hence $ \lambda: S\to \operatorname{End}_{\mathbb Z}(M\otimes_R N)$ is a ring homomorphism. Which means that $M\otimes_R N$ is a left $S$-module.

    Similarly, for fixed $t\in T$, define
    \[
        \gamma_t:M\times N\to M\otimes_R N,
        \qquad
        \gamma_t(m,n)=m\otimes nt.
    \]
    Again $\gamma_t$ is $R$-balanced, since for $r\in R$,
    \[
        \gamma_t(mr,n)=mr\otimes nt=m\otimes r(nt)=m\otimes (rn)t=\gamma_t(m,rn).
    \]
    Thus there exists a unique group homomorphism
    \[
        \rho_t:M\otimes_R N\to M\otimes_R N
    \]
    such that
    \[
        \rho_t(m\otimes n)=m\otimes nt.
    \]
    We write $\rho_t(x)=xt$. Hence $ \rho: T\to \operatorname{End}_{\mathbb Z}(M\otimes_R N)$ is a ring homomorphism. Which means that $M\otimes_R N$ is a right $T$-module.

    It remains to verify the bimodule axioms. Since elementary tensors generate $M\otimes_R N$, it suffices to check them on $m\otimes n$:
    \[
        (s_1+s_2)(m\otimes n)=((s_1+s_2)m)\otimes n
        =(s_1m)\otimes n+(s_2m)\otimes n
        =s_1(m\otimes n)+s_2(m\otimes n),
    \]
    \[
        (s_1s_2)(m\otimes n)=((s_1s_2)m)\otimes n
        =s_1((s_2m)\otimes n)=s_1\bigl(s_2(m\otimes n)\bigr),
    \]
    \[
        1_S(m\otimes n)=(1_Sm)\otimes n=m\otimes n,
    \]
    and similarly
    \[
        (m\otimes n)(t_1+t_2)=m\otimes n(t_1+t_2)
        =(m\otimes nt_1)+(m\otimes nt_2),
    \]
    \[
        ((m\otimes n)t_1)t_2
        =(m\otimes nt_1)t_2
        =m\otimes (nt_1)t_2
        =m\otimes n(t_1t_2)
        =(m\otimes n)(t_1t_2),
    \]
    \[
        (m\otimes n)1_T=m\otimes n.
    \]
    Finally,
    \[
        \bigl(s(m\otimes n)\bigr)t=((sm)\otimes n)t=(sm)\otimes nt
        =s(m\otimes nt)=s\bigl((m\otimes n)t\bigr).
    \]
    Therefore $M\otimes_R N$ is a left $S$-right $T$-bimodule.
\end{proof}

\begin{lemma}
    Let $M_R$,$M'_R$ be right $R$-modules, and let ${}_R N$,${}_R N'$ be left $R$-modules. Suppose
    \[
        f_1,f_2,f\in \Hom_R(M,M'),
        \qquad
        g_1,g_2,g\in \Hom_R(N,N').
    \]
    Then for each pair $(f,g)$ there exists a unique group homomorphism
    \[
        f\otimes g:M\otimes_R N\to M'\otimes_R N'
    \]
    such that
    \[
        (f\otimes g)(m\otimes n)=f(m)\otimes g(n)
        \qquad
        (m\in M,\ n\in N).
    \]
    Moreover,
    \begin{enumerate}
        \item $(f_1+f_2)\otimes g=(f_1\otimes g)+(f_2\otimes g)$;
        \item $f\otimes (g_1+g_2)=(f\otimes g_1)+(f\otimes g_2)$;
        \item $f\otimes 0=0$ and $0\otimes g=0$;
        \item $1_M\otimes 1_N=1_{M\otimes_R N}$.
    \end{enumerate}
\end{lemma}

\begin{proof}
    Define
    \[
        \beta:M\times N\to M'\otimes_R N',
        \qquad
        \beta(m,n)=f(m)\otimes g(n).
    \]
    We check that $\beta$ is $R$-balanced. It is additive in each variable because $f$ and $g$ are module homomorphisms. Also, for $r\in R$,
    \[
        \beta(mr,n)=f(mr)\otimes g(n)=f(m)r\otimes g(n)
        =f(m)\otimes rg(n)=f(m)\otimes g(rn)=\beta(m,rn).
    \]
    Hence, by the universal property of $M\otimes_R N$, there exists a unique group homomorphism
    \[
        f\otimes g:M\otimes_R N\to M'\otimes_R N'
    \]
    such that
    \[
        (f\otimes g)(m\otimes n)=f(m)\otimes g(n).
    \]

    Now each asserted identity is verified on elementary tensors, and therefore on all of $M\otimes_R N$.

    For $(1)$, for all $m\in M$ and $n\in N$,
    \[
        \bigl(((f_1+f_2)\otimes g)(m\otimes n)\bigr)
        =(f_1+f_2)(m)\otimes g(n)
        =f_1(m)\otimes g(n)+f_2(m)\otimes g(n),
    \]
    while
    \[
        \bigl((f_1\otimes g)+(f_2\otimes g)\bigr)(m\otimes n)
        =(f_1\otimes g)(m\otimes n)+(f_2\otimes g)(m\otimes n)
        =f_1(m)\otimes g(n)+f_2(m)\otimes g(n).
    \]
    Hence
    \[
        (f_1+f_2)\otimes g=(f_1\otimes g)+(f_2\otimes g).
    \]

    The proof of $(2)$ is analogous:
    \[
        (f\otimes (g_1+g_2))(m\otimes n)
        =f(m)\otimes (g_1+g_2)(n)
        =f(m)\otimes g_1(n)+f(m)\otimes g_2(n),
    \]
    which equals
    \[
        ((f\otimes g_1)+(f\otimes g_2))(m\otimes n).
    \]

    For $(3)$,
    \[
        (f\otimes 0)(m\otimes n)=f(m)\otimes 0=0,
        \qquad
        (0\otimes g)(m\otimes n)=0\otimes g(n)=0,
    \]
    so both maps are the zero homomorphism.

    For $(4)$,
    \[
        (1_M\otimes 1_N)(m\otimes n)=1_M(m)\otimes 1_N(n)=m\otimes n,
    \]
    hence $1_M\otimes 1_N$ is the identity on $M\otimes_R N$.
\end{proof}

\begin{lemma}
    Let $M_R$,$M'_R$,$M''_R$ be right $R$-modules, and let ${}_R N$,${}_R N'$,${}_R N''$ be left $R$-modules. Let
    \[
        f:M\to M',\qquad f':M'\to M'',
        \qquad
        g:N\to N',\qquad g':N'\to N''
    \]
    be $R$-module homomorphisms. Then
    \[
        (f'\otimes g')\circ (f\otimes g)=(f'\circ f)\otimes (g'\circ g).
    \]
\end{lemma}

\begin{proof}
    Again it suffices to check the equality on elementary tensors. For $m\in M$ and $n\in N$,
    \[
        \bigl((f'\otimes g')\circ (f\otimes g)\bigr)(m\otimes n)
        =(f'\otimes g')\bigl(f(m)\otimes g(n)\bigr)
        =f'(f(m))\otimes g'(g(n)),
    \]
    while
    \[
        \bigl((f'\circ f)\otimes (g'\circ g)\bigr)(m\otimes n)
        =(f'\circ f)(m)\otimes (g'\circ g)(n)
        =f'(f(m))\otimes g'(g(n)).
    \]
    Thus both homomorphisms agree on the generators $m\otimes n$ of $M\otimes_R N$, hence they are equal.
\end{proof}

\begin{definition}
    Let ${}_S U_R$ be a bimodule.
    Define
    \[
        U \otimes_R - : R\text{-}\mathrm{Mod} \longrightarrow S\text{-}\mathrm{Mod}
    \]
    by
    \[
        M \longmapsto U \otimes_R M,
        \qquad
        f \longmapsto \operatorname{id}_U \otimes f.
    \]
\end{definition}

Similarly, Let ${}_R U_S$ be a bimodule.
\[
    - \otimes_R U : \mathrm{Mod}\text{-}R \longrightarrow \mathrm{Mod}\text{-}S
\]
by
\[
    N \longmapsto N \otimes_R U,
    \qquad
    g \longmapsto g \otimes \operatorname{id}_U.
\]

\begin{proposition}
    Let ${}_S U_R$ be a bimodule. Then
    \[
        U \otimes_R - : R\text{-}\mathrm{Mod} \longrightarrow S\text{-}\mathrm{Mod}
    \]
    is a covariant functor.

    Similarly, if ${}_R U_S$ is a bimodule, then
    \[
        - \otimes_R U : \mathrm{Mod}\text{-}R \longrightarrow \mathrm{Mod}\text{-}S
    \]
    is a covariant functor.
\end{proposition}

\begin{proof}
    We prove the first statement. For each left $R$-module $M$, the tensor product
    \[
        U \otimes_R M
    \]
    is naturally a left $S$-module via
    \[
        s(u \otimes m)=(su)\otimes m
        \qquad
        (s\in S,\ u\in U,\ m\in M).
    \]
    For each $R$-module homomorphism $f:M\to N$, define
    \[
        U\otimes_R f: U\otimes_R M \to U\otimes_R N
    \]
    by
    \[
        (U\otimes_R f)(u\otimes m)=u\otimes f(m).
    \]
    Equivalently,
    \[
        U\otimes_R f=\operatorname{id}_U\otimes f.
    \]
    This is well-defined and is an $S$-module homomorphism.

    It remains to verify the functorial properties. For the identity map $\operatorname{id}_M:M\to M$, we have
    \[
        U\otimes_R \operatorname{id}_M=\operatorname{id}_U\otimes \operatorname{id}_M
        =\operatorname{id}_{U\otimes_R M}.
    \]
    If
    \[
        M \xrightarrow{f} N \xrightarrow{g} P
    \]
    are $R$-module homomorphisms, then for every $u\otimes m\in U\otimes_R M$,
    \[
        (U\otimes_R (g\circ f))(u\otimes m)
        = u\otimes (g\circ f)(m)
        = g(u\otimes f(m))
        = ((U\otimes_R g)\circ (U\otimes_R f))(u\otimes m).
    \]
    Hence
    \[
        U\otimes_R (g\circ f)=(U\otimes_R g)\circ (U\otimes_R f).
    \]
    Therefore $U\otimes_R -$ is a covariant functor from $R\text{-}\mathrm{Mod}$ to $S\text{-}\mathrm{Mod}$.

    The proof of the second statement is entirely similar. If ${}_R U_S$ is a bimodule and $N$ is a right $R$-module, then
    \[
        N\otimes_R U
    \]
    is naturally a right $S$-module via
    \[
        (n\otimes u)s=n\otimes (us).
    \]
    For a homomorphism $g:N\to N'$ in $\mathrm{Mod}\text{-}R$, define
    \[
        g\otimes_R U = g\otimes \operatorname{id}_U : N\otimes_R U \to N'\otimes_R U
    \]
    by
    \[
        (g\otimes_R U)(n\otimes u)=g(n)\otimes u.
    \]
    One checks immediately that identities and compositions are preserved. Thus
    \[
        -\otimes_R U:\mathrm{Mod}\text{-}R\to \mathrm{Mod}\text{-}S
    \]
    is also a covariant functor.
\end{proof}

\begin{theorem}
    Let ${}_S U_R$ be a bimodule. Then the functor
    \[
        U\otimes_R - \;:\; R\text{-Mod}\longrightarrow S\text{-Mod}
    \]
    is right exact. In other words, if
    \[
        N' \xrightarrow{\,f\,} N \xrightarrow{\,g\,} N'' \longrightarrow 0
    \]
    is an exact sequence of left $R$-modules, then
    \[
        U\otimes_R N' \xrightarrow{\,\mathrm{id}_U\otimes f\,}
        U\otimes_R N \xrightarrow{\,\mathrm{id}_U\otimes g\,}
        U\otimes_R N'' \longrightarrow 0
    \]
    is exact.
\end{theorem}

\begin{proof}
    Since
    \[
        N' \xrightarrow{\,f\,} N \xrightarrow{\,g\,} N'' \longrightarrow 0
    \]
    is exact, we have
    \[
        \ker g=\operatorname{Im} f
        \qquad\text{and}\qquad
        N''\cong N/\operatorname{Im} f.
    \]
    Let
    \[
        \pi:N\longrightarrow N/\operatorname{Im} f
    \]
    be the canonical projection. Then there exists an isomorphism
    \[
        \alpha:N/\operatorname{Im} f \longrightarrow N''
    \]
    such that
    \[
        g=\alpha\circ \pi.
    \]

    Now consider the map
    \[
        \Phi:U\otimes_R N \longrightarrow U\otimes_R (N/\operatorname{Im} f),
        \qquad
        u\otimes n \longmapsto u\otimes \pi(n).
    \]
    By the universal property of tensor products,$\Phi$ is a well-defined $S$-module homomorphism.

    We claim that
    \[
        \ker \Phi=\operatorname{Im}(\mathrm{id}_U\otimes f).
    \]

    First, since $\pi\circ f=0$, functoriality gives
    \[
        \Phi\circ (\mathrm{id}_U\otimes f)
        =
        \mathrm{id}_U\otimes (\pi\circ f)
        =
        0.
    \]
    Hence
    \[
        \operatorname{Im}(\mathrm{id}_U\otimes f)\subseteq \ker\Phi.
    \]

    To prove the reverse inclusion, define
    \[
        \overline{\Phi}:(U\otimes_R N)/\operatorname{Im}(\mathrm{id}_U\otimes f)
        \longrightarrow
        U\otimes_R (N/\operatorname{Im} f)
    \]
    by
    \[
        \overline{\Phi}\bigl((u\otimes n)+\operatorname{Im}(\mathrm{id}_U\otimes f)\bigr)
        =
        u\otimes \pi(n).
    \]
    This is well defined because $\Phi$ vanishes on $\operatorname{Im}(\mathrm{id}_U\otimes f)$.

    Next define
    \[
        \Psi:U\times (N/\operatorname{Im} f)\longrightarrow
        (U\otimes_R N)/\operatorname{Im}(\mathrm{id}_U\otimes f)
    \]
    by
    \[
        \Psi(u,\overline{n})
        =
        (u\otimes n)+\operatorname{Im}(\mathrm{id}_U\otimes f),
    \]
    where $\overline{n}=n+\operatorname{Im} f$.

    We check that $\Psi$ is well defined. If $\overline{n}=\overline{n'}$, then
    \[
        n-n'\in\operatorname{Im} f,
    \]
    so there exists $x\in N'$ such that
    \[
        n-n'=f(x).
    \]
    Hence
    \[
        (u\otimes n)-(u\otimes n')
        =
        u\otimes (n-n')
        =
        u\otimes f(x)
        =
        (\mathrm{id}_U\otimes f)(u\otimes x)
        \in\operatorname{Im}(\mathrm{id}_U\otimes f).
    \]
    Thus
    \[
        (u\otimes n)+\operatorname{Im}(\mathrm{id}_U\otimes f)
        =
        (u\otimes n')+\operatorname{Im}(\mathrm{id}_U\otimes f).
    \]
    So $\Psi$ is well defined. It is clearly $R$-balanced, hence induces an $S$-module homomorphism
    \[
        \overline{\Psi}:U\otimes_R (N/\operatorname{Im} f)
        \longrightarrow
        (U\otimes_R N)/\operatorname{Im}(\mathrm{id}_U\otimes f)
    \]
    such that
    \[
        \overline{\Psi}(u\otimes \overline{n})
        =
        (u\otimes n)+\operatorname{Im}(\mathrm{id}_U\otimes f).
    \]

    It is immediate that $\overline{\Phi}$ and $\overline{\Psi}$ are inverse to each other. Therefore
    \[
        (U\otimes_R N)/\operatorname{Im}(\mathrm{id}_U\otimes f)
        \cong
        U\otimes_R (N/\operatorname{Im} f).
    \]
    It follows that
    \[
        \ker\Phi=\operatorname{Im}(\mathrm{id}_U\otimes f).
    \]

    Finally, since $\alpha:N/\operatorname{Im} f\to N''$ is an isomorphism, the induced map
    \[
        \mathrm{id}_U\otimes \alpha:
        U\otimes_R (N/\operatorname{Im} f)
        \longrightarrow
        U\otimes_R N''
    \]
    is an isomorphism. Moreover,
    \[
        \mathrm{id}_U\otimes g
        =
        (\mathrm{id}_U\otimes \alpha)\circ \Phi.
    \]
    Hence
    \[
        \ker(\mathrm{id}_U\otimes g)=\ker\Phi=\operatorname{Im}(\mathrm{id}_U\otimes f),
    \]
    and $\mathrm{id}_U\otimes g$ is surjective because $\Phi$ is surjective and $\mathrm{id}_U\otimes \alpha$ is an isomorphism.

    Therefore
    \[
        U\otimes_R N' \xrightarrow{\,\mathrm{id}_U\otimes f\,}
        U\otimes_R N \xrightarrow{\,\mathrm{id}_U\otimes g\,}
        U\otimes_R N'' \longrightarrow 0
    \]
    is exact.
\end{proof}


\begin{theorem}
    Let ${}_RU_{S}$ be a left-$R$-right-$S$-bimodule. Then the functor
    \[
        -\otimes_R U \;:\; \text{Mod}-R\longrightarrow \text{Mod}-S
    \]
    is right exact.
\end{theorem}

